\section{没有快速色散波的快重联：反例具体推翻了什么}\label{sec:withoutfdw}
\citet{L216}在强导引场两场模型中分别研究$d_e>\rho_s$和$\rho_s>d_e$。前者没有通常用来解释快重联的快速色散分支。超电阻从$10^{-11}$增到$10^{-8}$时，两流体算例的归一化率仍在约0.025--0.029；对应单流体率则近于$\eta_H^{0.252}$。这是同一类耗散扫描中的实测差别，不是从X形图像直接判定快慢。

其关键新发现是\textbf{微观耗散层的上游场和长度会自调节}。当耗散层比$h=\max(d_e,\rho_s)$薄，上游$B_\delta$随厚度一起减小；某些区间内层长度也随厚度缩短。因此，对固定$B_\delta$和固定内层长度推导出的耗散依赖，不一定是实际算例采取的路径。第一章式\eqref{eq:two_layer_rate}把这一观察重建成两层匹配。它解释给定$B_h,w_v$时的弱耗散依赖，却没有自洽选出这两个外部量。

长文\citet{L217}进一步比较流体与PIC。在有快速色散波的一组参数中，两模型峰值率相差不到约10\%，但启动和饱和时刻不同；PIC早期出现的次级岛、喷流弯曲也没有在流体中逐一重现。在无快速色散波的\emph{小盒}算例中，没有次级岛仍能快速重联。因此，该组证据同时排除了“快速色散波必须存在”和“次级岛必须出现”这两种过强的必要性表述；并未排除它们在其他参数区影响速率。

\subsection{重做早期扫描：区分改变声尺度和改变电子尺度}
长文第VII节重做类似Rogers等的扫描：固定$d_e$而减小$\rho_s$时，速率下降可以出现；但固定$\rho_s$、在仍无快速色散波的区间增大$d_e$，速率也会上升。两组都改变了$h$相对于初始电流片宽的比例，因而改变启动等待时间及启动时还能供应多强的上游场。按当时测得的$B_h^2$归一化后，差别明显缩小。这里的结论是早期扫描不能单独隔离色散波的因果作用，并非“线性波对重联毫无作用”。\citep{Rogers2001,L217}

\subsection{尺寸扫描没有同时放大所有方向}
Harris系列把流出方向长度从基准的0.5倍变到8倍，基准$L_{x0}=100h$；入流方向盒宽与初始片厚保持不变。局地上游场取离中面约$4h$的位置，按入流穿越时间平均，所得流体与PIC率约0.09--0.17。长文还排除最后约四分之一时窗，以减少可用磁通耗尽的影响。这个结果支持所测试几何中的弱尺寸依赖，但不能替代全几何自相似放大，也不能证明任何边界下都为0.1。

学习时可先对照短文图1、3，再读长文图14--19：前者检查耗散层如何变化，后者检查尺寸、归一化和扫描路径。两篇包含相关算例，不能作为两次完全独立的统计验证。\citep{L216,L217}

\section{受迫重联：同一个“快”字包含不同驱动和结束方式}
\subsection{Taylor墙位移：快峰值后仍有漫长慢尾}
\citet{L057}研究最终停止移动的导体墙，给出线性和非线性、MHD和Hall的不同峰值标度。低电阻时，超电阻对X点电场的贡献会超过普通电阻；即使扫描中率对超电阻弱敏感，也不能把实际计算说成只含普通电阻的Hall模型。

任意导引场的四场扩展在$\beta_F=10^{-2}$至$10^2$范围、保持$d_\beta$相同时，得到相似的Hall突发阶段。保持的是$d_\beta=c_\beta d_i$，并非$d_i$固定。典型快峰值约为电阻参照的30倍，并有约8个Alfvén时间的延迟；\textbf{第一轮快爆发只把岛宽增长到最终值约四分之一}，随后仍有慢的Rutherford阶段。因此，“出现快速峰值”不等于“全部受迫磁通迅速完成重联”。\citep{L058}

两文将自身结果与近0.1的说法比较，但驱动、局地场和速率归一化不同。尤其后一文讨论中对上游磁场的线性标度与通常$B_uV_{Au}\propto B_u^2/\sqrt{n_u}$不同，不能只照搬其反驳措辞。更可靠的结论是：有限墙位移问题中的峰率和完成时间需要分别预测。

\subsection{体力持续压入：磁通堆积改变速率的分母}
\citet{L076}从较宽、撕裂稳定或边缘稳定的正弦片开始，用局域离子体力驱动入流。较弱驱动可建立平衡；较强驱动先积累上游磁场、压薄片，约到当地$0.7$--$0.8d_i$时出现快速阶段。其具体扫描中的临界驱动幅度在0.005--0.0075之间，这是选定盒子、驱动宽度和闭合的结果。

增加驱动宽度可提高\emph{绝对}电场，却因堆积的$B_u$更强而降低$B_uV_{Au}$归一化后的率。与自发薄片对照时，受迫算例的归一化出流和率可低约50\%。另一个值得单独记录的现象是：\textbf{磁岛宽度几乎饱和，岛内磁场却继续增加，故重联磁通仍增长}。用岛宽作为唯一通量代理会错判重联已经停止。

这篇2005年论文与2008年同题论文研究不同维数，不能合并成一个文件。2005年表I、正文算例交叉引用及部分图的驱动力数字有不一致；报告保留其主要可核实趋势，不为缺失输入自行补出精确对应。

\subsection{局域三维驱动：大部分磁通可以绕过重联层}
2008年研究把驱动同时沿第三方向局域化。二维中被驱动区域推来的上游磁通，超过70\%经过重联；相应三维例中约为10\%，其余可以沿第三方向偏转离开。三维局地电子电流还把磁结构向负$z$搬运，最大重联率不再位于驱动中心。持续显著重联的阶段也较短，典型比较约23对38个归一化时间。\citep{F717}

这里三维总磁通由场线追踪和面积积分求得，二维值则按单位第三维长度定义。作者将三维总率除以驱动全宽$2w_z=10d_i$作比较；这是一种有效长度归一化，不是一般三维重联率的唯一定义。典型对照的体力幅度也分别为二维0.15、三维0.2，并非仅改变维数的一对完全相同输入。结论应是\emph{局域第三维提供了磁通旁路，并改变局地层与整体驱动的联系}，而不是三维必然比二维快或慢一个固定倍数。

\section{陀螺诱导双重加速：先出现细尺度强流，再打开宏观排气}
\citet{F672}用强导引场四场陀螺流体，在大$\Delta'$、热离子和微观尺度远小于系统尺度时看到：线性增长之后先急剧加速，继而短暂减速或停滞，随后再次加速进入快重联。原文图6、8、9、11应连着看，分别追踪增长史、静电势、流场和排气开口；只看峰值会把两个过程合并。

第一轮加速与热离子有限回旋半径有关。在保持组合热尺度$\rho_\tau$相同、互换热离子／热电子贡献的对照中，热离子例出现更强的局地电势细结构，$E\times B$流可约高一个数量级；冷离子例缺少同样的第一轮加速。第一章已从准中性算子说明：热离子内区电势可直接跟随细密度结构，冷离子内区则经逆Laplacian平滑。这是与可测场量相连的机制对照。

第二轮加速随排气结构进一步开放，冷热离子均可出现。小尺度与宏观尺度分离越明显，两轮之间的减速越清楚。原文所画$\gamma_{\rm eff}=\dd\ln\delta\psi/\dd t$是\emph{对数磁通增长率}；它的尖峰不等于重联电场本身的尖峰，不能拿来与0.1直接比较。

表I将$d_e$从0.2减至0.025时，归一化峰值率由约0.283逐渐接近0.100；分母采用\emph{初始最大面内磁场}，不是随时间测得的局地上游场。作者也限定$d_e\Delta'\gtrsim1$：若固定$\Delta'$继续令$d_e\to0$，最终会进入小岛饱和区。因此该表不能证明严格理想极限仍有普适有限率。谱滤波、最细结构失去解析后的停止时刻和电子回旋效应缺失，同样限制了后期判断。

\subsection{哪些部分已经有动理学对照，哪些还没有}
\citet{E475}比较四场陀螺流体与EUTERPE陀螺动理学，先让两种代码各自对照其模型的线性本征解，再进行跨模型比较。这能把“离散是否解对自身方程”与“两个模型是否近似同一物理”分开。

在所扫线性波数中，增长率差别可达约20\%；减小平衡面内场后差别减弱，但没有自动降到零。非线性部分主要检验\emph{小}$\Delta'$：两模型都得到小岛饱和及其振荡，振荡频率很接近；饱和半宽差在较大驱动的$k_y=1.6$处约30\%，接近边缘稳定时明显减小。岛宽对温度相对不敏感，振荡频率却可随离子温度增加而近乎加倍。

论文文字称一个温比扫描的岛宽仅变约5\%，但图11显示的变化明显更大；据图估读，动理学约从0.040到0.049，流体约从0.046到0.053。这里只把它判为原文内部数量不一致，不把图上估读当作精确源数据。更关键的是，动理学非线性计算保留的模主要为$m=0,\pm1$，参数也在小岛区；\textbf{它没有检验前文大$\Delta'$的两次加速}。振荡可与粒子反弹频率比较，不代表有限矩模型恢复了完整俘获分布。

\section{加入Landau闭合：岛面积相近，增长时钟与温度结构已改变}
\citet{L244}构造含电子平行／垂直温度和离子垂直温度的六场模型，实际重联模拟则使用冷离子子系统。保留非局地电子热流后，线性增长下降；某些较高斜向模数的模式由不稳定变为稳定。闭合在谐振面局地为零，不妨碍它通过面外响应改变整个本征模。

非线性图像若用各自的$\gamma t$作横轴，岛形和最终面积可以相近，但实际时间已经不同。不能把这种重标度后的重合写成“Landau阻尼不改变重联速度”。温度、密度和电流的细结构也变平滑，所示电子平行温度峰值约降为无热流情况的一半。新发现是\textbf{宏观岛面积对闭合相对稳健，细结构和热响应却不稳健}。

其热流自由能损失可作为向未解析速度空间转移的代理，但模型没有复现所引用动理学研究中磁能释放与电子加热之间的明显时间延迟；作者提出闭合可能耗散过强。另一个量化限制是，约$0.002H_0$的数值能量误差已占本次释放磁能约20\%。故能量分配的精确比例没有像磁岛形态那样受到充分验证。

这些算例只含单一螺旋性，可以用$y^*=y+(k_z/k_y)z$约成二维；它们验证斜向结构与沿场闭合，却不能等同于多个斜向模相互作用的完整三维湍流。原文图4--7和能量段应一起读，避免只选最相似的岛图而忽略时间、闭合耗散与能量误差。
